\chapter{Plant Based Diet}

\medskip
\noindent\textit{\textbf{Under review.} This chapter is an AI-assisted educational draft; it is not a substitute for professional medical advice.}

\section*{Definition and Overview}
a plant-based diet emphasises vegetables, fruit, legumes, whole grains, nuts, seeds, and plant-derived foods. It may include or exclude animal products depending on the person's chosen pattern.

\section*{Clinical Features}
there are no disease-specific symptoms; poorly planned diets may cause inadequate protein, vitamin B12, iron, calcium, iodine, vitamin D, or omega-3 intake.

\section*{When to Seek Medical Care}
Seek medical review for new, persistent, worsening, or function-limiting symptoms. Call 000 for severe breathing difficulty, collapse, sudden neurological change, severe chest or abdominal pain, or any rapidly worsening illness.

\section*{Diagnosis and Investigations}
review dietary pattern, growth or weight, medical conditions, medicines, and targeted nutrient tests when deficiency is suspected.

\section*{Management}
eat varied protein and fortified foods, obtain reliable B12 sources, and seek dietitian advice during pregnancy, childhood, illness, or restrictive diets.

\section*{Complications}
Complications depend on the underlying cause and may include progression, effects on other organs, treatment adverse effects, or reduced quality of life.

\section*{Prognosis and Self-Care}
Outcome depends on the cause, severity, complications, and response to treatment. Follow the care plan, keep appointments, avoid known triggers, and seek help if symptoms worsen.

\section*{Expanded clinical context}
This chapter addresses plant based diet, a topic that should be interpreted in the context of age, baseline health, medicines, comorbidities, goals, and access to care. It supports discussion with a qualified clinician and does not establish a diagnosis.

\section*{Assessment and decision points}
Review onset, duration, triggers, relevant exposures, physical findings, associated symptoms, functional impact, and immediate safety. Record the main concern, time course, risk factors, red flags, and the person’s questions. Use targeted investigations when their result is likely to change management.

\section*{Management and follow-up}
Care may include education, practical support, treatment of contributing conditions, medication or a procedure when indicated, and a shared follow-up plan. Explain expected improvement, possible adverse effects, and when reassessment is needed. Avoid starting, stopping, or sharing prescription medicines without professional advice.

\section*{Safety-netting}
Seek urgent help for severe or rapidly worsening symptoms, collapse, breathing difficulty, severe pain, confusion, dehydration, uncontrolled bleeding, or any immediate danger.

\section*{Practical prevention}
Use plain-language information, supportive routines, risk reduction, and professional review when symptoms persist, recur, or impair daily life. Keep written instructions and confirm how to access help outside routine appointments.

\section*{Limitations}
This educational expansion should be checked against current local guidance and authoritative topic-specific sources.
\clearpage
